\documentclass[11pt,a4paper]{article}

\usepackage[utf8]{inputenc}
\usepackage[T1]{fontenc}
\usepackage{lmodern}
\usepackage{amsmath,amssymb,amsthm}
\usepackage{booktabs}
\usepackage{enumitem}
\usepackage{hyperref}
\usepackage[margin=2.5cm]{geometry}
\usepackage{xcolor}

\definecolor{codeblue}{RGB}{30,90,156}

\newtheorem{theorem}{Theorem}[section]
\newtheorem{proposition}[theorem]{Proposition}
\newtheorem{corollary}[theorem]{Corollary}
\newtheorem{lemma}[theorem]{Lemma}

\theoremstyle{definition}
\newtheorem{definition}[theorem]{Definition}

\theoremstyle{remark}
\newtheorem*{remark}{Remark}

\hypersetup{
  colorlinks=true,
  linkcolor=codeblue,
  citecolor=codeblue,
  urlcolor=codeblue,
}

\newcommand{\Qpos}{\mathbb{Q}_{>0}}
\newcommand{\Rpos}{\mathbb{R}_{>0}}
\newcommand{\costL}[1]{F_{#1}}

\title{%
  The Cost Form Is Forced, the Unit Is a Gauge:\\
  How Far Distinction Determines the Recognition Cost}
\author{%
  Jonathan Washburn\\
  Recognition Physics Institute, Austin, TX, USA\\
  \texttt{jon@recognitionphysics.org}}
\date{May 2026}

\begin{document}

\maketitle

\begin{abstract}\sloppy
The Recognition Science program rests on a single cost function, the canonical
reciprocal cost $J(x) = \tfrac{1}{2}(x + x^{-1}) - 1$, and on the claim that this
cost is forced rather than chosen. We determine exactly how far that claim
holds. We lead with the negative, because the negative is where the honest
account begins. On the field of ratios that the primitive act of distinction
generates, the cost is \emph{not} forced. The multiplicative group of positive
ratios is free abelian on the primes, the reciprocal-symmetric solutions of the
recognition composition law on that group are parameterized by its characters,
and a character is fixed by freely assigning a value to each prime. The space of
admissible costs is therefore an infinite product, one independent factor per
prime, and the canonical cost is a single point in it. We exhibit the freedom
concretely: for every prime there is an admissible cost that agrees with the
canonical one on all other primes and inverts it on that one. Any uniqueness
statement that recovers the canonical cost on the ratios must first fix the cost
on a generating set of primes, which is to assume agreement with the answer, not
to derive it. We then show what the continuous completion adds. A continuous
character of the positive reals is a power map, so on the completion the four
algebraic laws cut the infinite character freedom down to a single one-parameter
family, $\costL{\lambda}(x) = \tfrac{1}{2}(x^{\lambda} + x^{-\lambda}) - 1$ for
$\lambda > 0$, and that family is exactly the solution set, not merely contained
in it. The family is the orbit of the canonical cost under the automorphism
group of the positive reals under multiplication, and the action is free and
transitive, so the family is a torsor and the residual freedom is exactly one
positive real, pinned by agreement at a single ratio. A unit calibration selects
the canonical cost within the family. The honest conclusion is a stratification:
distinction and the algebraic laws force the cost \emph{form}, the gauge orbit of
$J$; one calibration datum fixes the gauge; and the value of that datum is a
multiplicative-automorphism gauge that distinction does not, by itself, fix. We
close by naming the one open question that would upgrade the result, namely
whether the calibration unit is itself derivable from the cost of a single
distinction act.
\end{abstract}

\tableofcontents

\section{The joint, and the claim this paper defends}

A program that derives physics from a single primitive has one load-bearing
joint, and in the Recognition Science program it is the cost function. The
program assigns to each comparison of two magnitudes a cost, and almost
everything downstream, the appearance of the golden ratio, the discrete time
cadence, the spatial dimension, the numerical values of physical constants,
flows from the precise form of that cost. The cost is the canonical reciprocal
cost
\[
  J(x) = \tfrac{1}{2}\bigl(x + x^{-1}\bigr) - 1 ,
\]
and the program's strong claim is that this cost is not one choice among many but
the forced consequence of a few structural requirements. If that claim is true,
the cost layer adds nothing arbitrary. If it is false, the cost layer hides a
choice, and the program owes an account of what the choice is and how large it
is.

This paper gives that account, and the account is a stratification rather than a
clean yes or no. We state the conclusion at the outset so that nothing in what
follows can be mistaken for a stronger claim than we make. On the discrete field
of ratios that distinction generates, the cost is not forced: the admissible
costs form an infinite-dimensional family and the canonical cost is one member.
On the continuous completion of that field, the same algebraic laws, now joined
by continuity, force a single one-parameter family, and that family is precisely
the orbit of the canonical cost under a natural group of symmetries. Within the
family a single calibration datum selects the canonical cost. So distinction
forces the cost \emph{form}, and one unit of scale, supplied by a single
measurement, selects the canonical representative. The value of that unit is a
gauge that distinction does not fix on its own.

We organize the paper negative first, and the ordering is deliberate. The most
common way a forcing claim of this kind fails is by circularity: the argument
that appears to force the cost turns out to have assumed the answer on a set rich
enough to determine it. The cleanest defense against that failure is to confront
the negative before the positive, to show exactly where and why the cost is
\emph{not} determined, and only then to identify the minimal extra ingredient
that determines it. A reader who sees the obstruction first cannot later mistake
a hidden assumption for a derivation. This paper therefore proves the negative in
full, names the circularity trap explicitly, and then builds the positive on the
completion in a way that keeps the one remaining choice in plain view.

\section{The cost and its laws}

We recall the structural requirements that a cost is asked to satisfy. They are
stated for a function on the positive reals, but the first three are purely
algebraic and make sense on any field of values, in particular on the rational
field that distinction produces.

\begin{definition}[The cost laws]\label{def:laws}
Let $K$ be an ordered field and let $K_{>0}$ be its positive part under
multiplication. A function $F \colon K_{>0} \to K$ is said to satisfy:
\begin{enumerate}[leftmargin=2.6em, label=\textup{(L\arabic*)}]
  \item \textbf{reciprocal symmetry}, if $F(x) = F(x^{-1})$ for all $x$;
  \item \textbf{normalization}, if $F(1) = 0$;
  \item \textbf{the composition law}, if for all $u, v$,
    \[
      F(uv) + F(u/v) = 2F(u)F(v) + 2F(u) + 2F(v) ;
    \]
  \item \textbf{continuity}, when $K = \mathbb{R}$, if $F$ is continuous on
    $\Rpos$;
  \item \textbf{unit calibration}, when $K = \mathbb{R}$ and $F$ is twice
    differentiable in the logarithmic coordinate, if the second derivative of
    $t \mapsto F(e^{t})$ at $t = 0$ equals $1$.
\end{enumerate}
We call $F$ a \emph{cost form} if it satisfies (L1) through (L3), and a
\emph{calibrated cost} if in addition it satisfies (L4) and (L5).
\end{definition}

The first three laws are the algebra of comparison. Reciprocal symmetry says
that comparing $a$ to $b$ costs what comparing $b$ to $a$ costs. Normalization
says that comparing a thing to itself costs nothing. The composition law is the
nontrivial one: it is the two-branch identity that a recognition cost must obey
when comparisons are composed, and it is itself forced, among symmetric
right-affine combiners with natural boundary conditions, as a separate result
that we take as given here and cite to the companion functional-equation
work~\cite{WashburnFE2026}. The canonical cost $J$ satisfies all five laws. The
question of this paper is whether the laws, in turn, force $J$, and the answer
depends sharply on whether the domain is the discrete rational field or its
continuous completion.

It is worth isolating the role of calibration before going further, because it
is the hinge of the whole analysis. Calibration (L5) is a statement about a
second derivative at a single point. A second derivative is a local analytic
quantity: it sees the behavior of the function in an arbitrarily small
neighborhood of the unit. On a domain riddled with gaps, such as the rationals,
there is no neighborhood to see, and calibration has no meaning. This is the
first sign that the discrete and continuous cases will diverge, and it tells us
in advance where the divergence will come from.

\section{The negative: on the ratios, the cost is not forced}\label{sec:negative}

We prove the negative result first and in its strongest form: not as a list of
counterexamples but as a structural description of the entire space of admissible
costs on the rationals. The description shows at once that the space is large,
that the canonical cost is one unremarkable point in it, and that no argument
internal to the discrete laws can single the canonical cost out.

\subsection{The positive rationals are free abelian on the primes}

\begin{lemma}[Structure of the rational multiplicative group]\label{lem:free}
The group $\Qpos$ of positive rationals under multiplication is free abelian on
the set of primes. Concretely, the map sending a positive rational to its tuple
of prime exponents,
\[
  q \;\longmapsto\; \bigl(v_p(q)\bigr)_{p \text{ prime}} ,
\]
is an isomorphism from $\Qpos$ onto the direct sum $\bigoplus_{p} \mathbb{Z}$, the
tuples of integers that are zero in all but finitely many coordinates.
\end{lemma}

\begin{proof}
This is the fundamental theorem of arithmetic in multiplicative form. Every
positive rational has a unique expression $q = \prod_p p^{v_p(q)}$ with integer
exponents $v_p(q)$, all but finitely many zero, and the exponent of each prime
adds under multiplication of rationals. Uniqueness of prime factorization makes
the exponent map a bijection onto the finitely supported integer tuples, and the
additivity of exponents makes it a group homomorphism. Hence it is an
isomorphism, and the primes are a free basis.
\end{proof}

The content of Lemma~\ref{lem:free} for the cost question is that the primes are
\emph{independent} coordinates. A function defined on $\Qpos$ that respects the
multiplicative structure is free to behave on each prime axis without regard to
the others, because the axes do not interact. This independence is what makes the
space of admissible costs large.

\subsection{Cost forms are characters}

We now solve the cost laws on the rational group exactly. The solution is the
classical general solution of the composition law on an abelian group, and it
identifies cost forms with characters.

\begin{definition}[Character]\label{def:char}
A \emph{character} of $\Qpos$ into an ordered field $K$ is a group homomorphism
$\chi \colon \Qpos \to K_{>0}$, that is, a map with $\chi(uv) = \chi(u)\chi(v)$
and $\chi(1) = 1$.
\end{definition}

\begin{proposition}[Cost forms are character costs]\label{prop:charcost}
Let $K$ be an ordered field in which $2$ is invertible. A function
$F \colon \Qpos \to K$ is a cost form, that is satisfies reciprocal symmetry,
normalization, and the composition law, if and only if there is a character
$\chi$ of $\Qpos$ into $K$ with
\[
  F(u) = \tfrac{1}{2}\bigl(\chi(u) + \chi(u)^{-1}\bigr) - 1
  \qquad\text{for all } u \in \Qpos .
\]
\end{proposition}

\begin{proof}
Suppose first that $F$ has the stated form for a character $\chi$. Normalization
holds because $\chi(1) = 1$ gives $F(1) = \tfrac{1}{2}(1 + 1) - 1 = 0$.
Reciprocal symmetry holds because $\chi(u^{-1}) = \chi(u)^{-1}$, so the two terms
$\chi(u)$ and $\chi(u)^{-1}$ simply exchange. For the composition law, write
$a = \chi(u)$ and $b = \chi(v)$, so $\chi(uv) = ab$ and $\chi(u/v) = a/b$, and
set $H(w) = F(w) + 1 = \tfrac{1}{2}(\chi(w) + \chi(w)^{-1})$. Then
\[
  H(uv) + H(u/v)
  = \tfrac{1}{2}\Bigl(ab + \tfrac{1}{ab} + \tfrac{a}{b} + \tfrac{b}{a}\Bigr)
  = \tfrac{1}{2}\Bigl(a + \tfrac{1}{a}\Bigr)\Bigl(b + \tfrac{1}{b}\Bigr)
  = 2 H(u) H(v) ,
\]
and substituting $H = F + 1$ and expanding gives exactly the composition law for
$F$.

Conversely, suppose $F$ is a cost form, and set $H = F + 1$, so $H(1) = 1$,
$H(u) = H(u^{-1})$, and $H(uv) + H(u/v) = 2H(u)H(v)$. This last is d'Alembert's
functional equation on the abelian group $\Qpos$. Its general solution on an
abelian group, into a field where $2$ is invertible, is
$H(u) = \tfrac{1}{2}(\chi(u) + \chi(u)^{-1})$ for a homomorphism $\chi$ into the
multiplicative group, a classical fact about d'Alembert's equation
\cite{Aczel1966}. Concretely, one recovers the character by solving the quadratic
$\chi(u)^2 - 2H(u)\chi(u) + 1 = 0$ consistently across the group; the two roots
are $\chi(u)$ and $\chi(u)^{-1}$, and the composition law is exactly the
condition that a consistent multiplicative choice of root exists. Setting
$F = H - 1$ returns the stated form.
\end{proof}

Proposition~\ref{prop:charcost} reduces the entire cost question on the rationals
to a question about characters, and the canonical cost is the character cost of
the simplest possible character.

\begin{corollary}[The canonical cost is the identity character]\label{cor:Jchar}
The canonical reciprocal cost $J$ on $\Qpos$ is the character cost of the
inclusion character $\iota(u) = u$. That is, $J(u) = \tfrac{1}{2}(u + u^{-1}) - 1$
is $F$ with $\chi = \iota$.
\end{corollary}

\subsection{The freedom, and the canonical cost as one point in it}

We can now read off the size of the admissible family. A character of $\Qpos$ is,
by Lemma~\ref{lem:free}, a homomorphism out of a free abelian group, and such a
homomorphism is determined by, and may be freely prescribed by, its values on the
free generators, the primes.

\begin{theorem}[The cost is not forced on the rationals]\label{thm:negative}
The cost forms on $\Qpos$ with values in $\Qpos$ are in bijection with the
characters of $\Qpos$, and the characters are in bijection with the choices of an
image $\chi(p) \in \Qpos$ for each prime $p$, freely and independently. The space
of cost forms is therefore the infinite product $\prod_{p} \Qpos$, one
independent factor per prime. The canonical cost $J$ is the single point given by
$\chi(p) = p$ for every prime. No subset of the cost laws (L1) through (L3)
singles out that point.
\end{theorem}

\begin{proof}
By Proposition~\ref{prop:charcost} the cost forms correspond to characters, and by
Lemma~\ref{lem:free} a character is a homomorphism from the free abelian group
$\bigoplus_p \mathbb{Z}$ to $\Qpos$. A homomorphism out of a free abelian group is
determined by its values on a basis, and any assignment of basis values extends
to a unique homomorphism. The basis is the set of primes, so a character is
exactly a free choice of $\chi(p) \in \Qpos$ for each prime, and distinct choices
give distinct characters, hence distinct cost forms. The canonical cost is the
character $\chi = \iota$, that is $\chi(p) = p$. Since the laws (L1) through (L3)
are precisely the conditions that make $F$ a character cost, and every character
yields one, the laws hold for the entire product and cannot distinguish one point
from another.
\end{proof}

The abstract statement becomes vivid in a single explicit family of
counterexamples, one for each prime, each of which satisfies every discrete law
and differs from the canonical cost.

\begin{proposition}[Per-prime orientation is free]\label{prop:twist}
Fix a prime $p_0$. Define $\chi_{p_0}$ on generators by $\chi_{p_0}(p_0) = p_0^{-1}$
and $\chi_{p_0}(p) = p$ for every prime $p \neq p_0$, and extend multiplicatively.
Then $\chi_{p_0}$ is a character, the associated cost form $F_{\chi_{p_0}}$
satisfies reciprocal symmetry, normalization, and the composition law, and it
differs from the canonical cost $J$. Explicitly $\chi_{p_0}(q) = q \cdot p_0^{-2 v_{p_0}(q)}$,
so $F_{\chi_{p_0}}$ inverts the orientation of the $p_0$ axis and leaves every
other axis exactly as in $J$.
\end{proposition}

\begin{proof}
The prescription assigns a value in $\Qpos$ to each free generator, so it extends
to a unique character by Lemma~\ref{lem:free}, and
Proposition~\ref{prop:charcost} makes the associated cost form satisfy the three
laws. The closed form follows by evaluating on $q = \prod_p p^{v_p(q)}$:
$\chi_{p_0}(q) = p_0^{-v_{p_0}(q)} \prod_{p \neq p_0} p^{v_p(q)}
= q \cdot p_0^{-2 v_{p_0}(q)}$. It differs from $J$ already at $q = p_0$, where
$\chi_{p_0}(p_0) = p_0^{-1} \neq p_0 = \iota(p_0)$, so the two character costs
disagree at $p_0$ unless $p_0 = p_0^{-1}$, which never holds for a prime.
\end{proof}

\subsection{Why a discrete uniqueness theorem would be circular}

The space of cost forms on the rationals being an infinite product, one can still
prove a uniqueness statement of the following shape: if a cost form agrees with
$J$ on enough primes, then it equals $J$ everywhere. This is true, and it is
worth being completely clear that it is not a forcing result.

\begin{proposition}[Conditional, and circular as a forcing claim]\label{prop:circular}
A cost form $F$ on $\Qpos$ equals $J$ if and only if its character $\chi$ agrees
with the inclusion on every prime, that is $\chi(p) = p$ for all $p$. In
particular, fixing the cost on the full generating set of primes determines it
everywhere.
\end{proposition}

\begin{proof}
A homomorphism out of a free abelian group is determined by its values on the
basis, so two characters agree everywhere if and only if they agree on every
prime. Apply this to $\chi$ and $\iota$.
\end{proof}

Proposition~\ref{prop:circular} says exactly what it says and no more. It tells
us that the primes generate, so agreement on the primes propagates. It does not
tell us that distinction selects the inclusion character rather than any of the
infinitely many others. A uniqueness theorem whose hypothesis is ``the cost
already equals $J$ on a generating set'' has assumed the answer on a set large
enough to pin it down, and then observed that the answer propagates. That is a
fact about generation, not a derivation of $J$ from distinction. We state this
plainly because the temptation to present such a conditional as evidence of
forcing is real, and a careful reader sees through it at once. The honest reading
of the discrete layer is the one given by Theorem~\ref{thm:negative}: the cost is
not forced, and the freedom is an entire character group.

\section{Why the discrete carrier cannot do better}\label{sec:obstruction}

Before turning to the completion, we record why the negative of
Section~\ref{sec:negative} is structural rather than a temporary gap. One might
hope that some further discrete condition, short of completing the field, could
cut the character freedom down. It cannot, and the reason is the same reason the
calibration law had no meaning on the rationals.

The character freedom is killed only by continuity. A character of $\Qpos$ is an
arbitrary assignment of values on the primes; to constrain it one needs a
condition that links the value at one rational to the values at nearby rationals.
But on the rationals there is no useful notion of ``nearby'' in the order
topology that the cost laws can exploit, because the rationals are totally
disconnected: between any two of them lie gaps, and a constraint imposed near one
point does not propagate across a gap to a distant point. The per-prime
independence of Lemma~\ref{lem:free} is the algebraic face of the same fact. The
prime axes are not merely independent as coordinates; they cannot be coupled by
any order- or limit-based condition available on the discrete field, because such
conditions require a connected domain to have force.

This is the precise sense in which the cost question and the continuum question
are the same question seen twice. The cost is underdetermined on the rationals
for exactly the reason the rationals are not complete: the domain is full of
gaps, and gaps insulate constraints. To force the cost one must first remove the
gaps, and removing the gaps is completing the field to the continuum. As shown in
the companion analysis of the continuum~\cite{WashburnContinuum2026}, that
completion is not itself a consequence of distinction but a separate posit. So the
cost can be forced only on a domain that is itself an addition to distinction, and
the logical order is fixed: domain first, then the cost on it.

\section{The completion: continuity collapses the freedom}\label{sec:completion}

On the continuous completion of the rationals, the positive real line, the cost
laws acquire the two analytic conditions, continuity and calibration, that have
no meaning on the gapped field. We now show that continuity alone collapses the
infinite character freedom to a single one-parameter family, and that calibration
then selects the canonical cost within it.

The mechanism is that continuity restricts characters drastically. On the
rationals a character could send each prime anywhere; on the reals a continuous
character is forced to be a power map.

\begin{lemma}[Continuous characters of the reals are power maps]\label{lem:power}
A continuous homomorphism $\chi \colon \Rpos \to \Rpos$ is of the form
$\chi(x) = x^{\lambda}$ for a unique real $\lambda$.
\end{lemma}

\begin{proof}
Transport to logarithmic coordinates by $g(t) = \ln \chi(e^{t})$. Then $g$ is
continuous and additive, $g(s + t) = g(s) + g(t)$, since $\chi$ is multiplicative.
The continuous solutions of Cauchy's additive equation are the linear maps
$g(t) = \lambda t$ for a unique real $\lambda$. Hence
$\ln \chi(e^{t}) = \lambda t$, that is $\chi(x) = x^{\lambda}$, and $\lambda$ is
unique because distinct slopes give distinct maps.
\end{proof}

\begin{theorem}[The four laws admit exactly the one-parameter family]\label{thm:family}
Define, for $\lambda > 0$,
\[
  \costL{\lambda}(x) = \tfrac{1}{2}\bigl(x^{\lambda} + x^{-\lambda}\bigr) - 1 ,
  \qquad x \in \Rpos .
\]
Each $\costL{\lambda}$ is a continuous cost form, that is satisfies (L1) through
(L4). Conversely, every continuous cost form on $\Rpos$ whose logarithmic profile
has positive curvature at the unit equals $\costL{\lambda}$ for exactly one
$\lambda > 0$. The continuous cost forms are therefore precisely the family
$\{\costL{\lambda} : \lambda > 0\}$, together with the constant zero cost in the
degenerate flat case.
\end{theorem}

\begin{proof}
That each $\costL{\lambda}$ is a cost form is the character-cost computation of
Proposition~\ref{prop:charcost} applied to the continuous character
$\chi(x) = x^{\lambda}$, and continuity is clear. For the converse, let $F$ be a
continuous cost form. By the same proposition, applied now over the reals, $F$ is
the character cost of a continuous character $\chi$, and by Lemma~\ref{lem:power}
that character is $\chi(x) = x^{\lambda}$ for a unique real $\lambda$. Then
$F = \costL{\lambda}$. Reciprocal symmetry identifies $\lambda$ with $-\lambda$,
since $\costL{\lambda} = \costL{-\lambda}$, so we may take $\lambda \geq 0$; the
value $\lambda = 0$ gives the flat cost $F \equiv 0$, which has zero curvature,
and every $\lambda > 0$ gives a profile $t \mapsto \cosh(\lambda t) - 1$ with
positive curvature at the unit. Hence the nondegenerate continuous cost forms are
exactly $\{\costL{\lambda} : \lambda > 0\}$, parameterized injectively by
$\lambda > 0$.
\end{proof}

The classification of Theorem~\ref{thm:family} is the point at which the program's
continuous uniqueness result enters, and we have stated it so as to make explicit
what that result does and does not deliver. It delivers a one-parameter family. It
does not deliver a single cost, because calibration has not yet been imposed. The
single calibrated cost is recovered in Section~\ref{sec:gauge}, but only after we
understand the family as a whole, which is itself the structurally important
object.

\section{Monotonicity replaces continuity: the cost side needs no continuum}\label{sec:monotone}

The classification of Theorem~\ref{thm:family} was stated with continuity because
continuity is the hypothesis under which Cauchy's additive equation has only
linear solutions. Continuity is more than the argument needs. The additive
solutions can be pinned by an order condition instead, and the order condition is
available on fields far smaller than the continuum.

\begin{lemma}[Monotone additive maps are linear]\label{lem:monotone-additive}
Let $g \colon \mathbb{R} \to \mathbb{R}$ satisfy $g(s+t) = g(s) + g(t)$ for all
$s, t$, and suppose $g$ is monotone on some interval. Then $g(s) = c\,s$ for a
constant $c$.
\end{lemma}

This is classical (Darboux; see also Acz\'el). A monotone additive function is
bounded on an interval, and a bounded additive function is linear. Monotonicity
supplies the boundedness that continuity also supplies, but it is the weaker
requirement: every continuous additive map is monotone on the relevant rays,
while a monotone additive map need not be assumed continuous in advance.

Replacing continuity by monotonicity in Lemma~\ref{lem:power} leaves its
conclusion intact. A monotone multiplicative endomorphism of $\Rpos$ is a power
map. Theorem~\ref{thm:family} then holds verbatim with ``continuous'' replaced by
``order preserving'': the nondegenerate \emph{monotone} cost forms are exactly the
gauge orbit $\{\costL{\lambda} : \lambda > 0\}$.

The gain is not cosmetic. Continuity is a statement about limits, and limits
require the gaps of the rational line to be filled, which is exactly the move to
the uncountable continuum. Monotonicity is a statement about order, and order is
already present on the real-algebraic closure of the distinction-native rationals.
That closure is the field in which every distinction-posable polynomial comparison
resolves. It contains $\sqrt{2}$ and $\sqrt{5}$, hence the golden ratio
$\varphi = (1 + \sqrt{5})/2$, and it is countable. On that countable, real-closed,
ordered field the monotone cost forms are already exactly the gauge orbit. The
cost form is forced without ever invoking the continuum.

This sharpens the division of labor between the negative and positive results. The
companion analysis of the continuum~\cite{WashburnContinuum2026} argued that
completing the rationals to the real line is a separate posit, not a consequence
of distinction. The present result removes the only place where the cost argument
seemed to need that posit. Continuity was the apparent dependency; monotonicity
discharges it. Nothing of the continuum survives on the cost side. The genuinely
transcendental physical constants that the downstream theory eventually names,
which use $\pi$ and the exponential, are not algebraic and so do not lie in the
algebraic closure; but they too live in a countable field, obtained by adjoining
those specific constant values to the rationals, and that field is again a proper
subset of the continuum. The continuum is at most a scaffold for defining the
transcendental functions, never the home of the answers.

The residual freedom is visible as a single analytic invariant. In logarithmic
coordinates each member of the orbit is $t \mapsto \cosh(\lambda t) - 1$, whose
curvature at the unit, the second derivative at $t = 0$, equals $\lambda^{2}$.
Distinct members have distinct curvature, and the canonical cost $J$ is the unique
member of curvature one. Fixing the gauge is fixing that curvature, and nothing in
the discrete distinction data, which lives before any second derivative exists,
selects the value one. The unit is a calibration, not a forced constant.

\section{The family is the gauge orbit of the canonical cost}\label{sec:gauge}

The one-parameter family of Theorem~\ref{thm:family} is not an unstructured set
of solutions. It is a single orbit of the canonical cost under a natural group
action, and recognizing it as an orbit is what turns ``the cost is forced up to
one parameter'' into a precise statement about the kind of freedom that remains.

\begin{definition}[The gauge group]\label{def:gauge}
The \emph{gauge group} is the group $G$ of continuous automorphisms of $\Rpos$
under multiplication. By Lemma~\ref{lem:power} an element of $G$ is a power map
$\sigma_{\mu} \colon x \mapsto x^{\mu}$ with $\mu \neq 0$, and composition
multiplies exponents, $\sigma_{\mu} \circ \sigma_{\nu} = \sigma_{\mu\nu}$. The
positive part $G_{+} = \{\sigma_{\mu} : \mu > 0\}$ is the connected component of
the identity, isomorphic to $(\Rpos, \cdot)$ and to $(\mathbb{R}, +)$ through the
exponent.
\end{definition}

\begin{theorem}[The family is a torsor under the gauge group]\label{thm:torsor}
The gauge group acts on cost forms by precomposition,
$(\sigma_{\mu} \cdot F)(x) = F(x^{\mu})$. Under this action:
\begin{enumerate}[leftmargin=2em, label=\textup{(\alph*)}]
  \item the canonical cost is carried to the family, $\sigma_{\lambda} \cdot J = \costL{\lambda}$,
    so every member of the family is a gauge image of $J$;
  \item the action of $G_{+}$ on the family is transitive: any member reaches any
    other, since $\costL{\lambda} = \sigma_{\lambda/\nu} \cdot \costL{\nu}$;
  \item the action is free: $\sigma_{\mu} \cdot \costL{\nu} = \costL{\nu}$ with
    $\mu, \nu > 0$ forces $\mu = 1$, equivalently distinct exponents give distinct
    costs.
\end{enumerate}
Consequently the family $\{\costL{\lambda} : \lambda > 0\}$ is a principal
homogeneous space, a torsor, under $G_{+} \cong (\mathbb{R}, +)$. The residual
freedom in the cost is exactly one real degree of freedom, the gauge coordinate.
\end{theorem}

\begin{proof}
For (a), precomposing $J$ with $\sigma_{\lambda}$ gives
$J(x^{\lambda}) = \tfrac{1}{2}(x^{\lambda} + x^{-\lambda}) - 1 = \costL{\lambda}(x)$.
For (b), $\sigma_{\lambda/\nu} \cdot \costL{\nu}(x) = \costL{\nu}(x^{\lambda/\nu})
= \tfrac{1}{2}(x^{\lambda} + x^{-\lambda}) - 1 = \costL{\lambda}(x)$, so any
$\costL{\nu}$ reaches any $\costL{\lambda}$. For (c), suppose
$\costL{\mu\nu} = \sigma_{\mu} \cdot \costL{\nu} = \costL{\nu}$ with $\mu, \nu > 0$.
Evaluating at any $x \neq 1$, the function $\lambda \mapsto \costL{\lambda}(x)
= \cosh(\lambda \ln x) - 1$ is strictly increasing in $\lambda > 0$ for fixed
$x \neq 1$, since $\cosh$ is strictly increasing on the positive axis and
$|\ln x| > 0$. Hence $\costL{\mu\nu}(x) = \costL{\nu}(x)$ forces $\mu\nu = \nu$,
that is $\mu = 1$. A free transitive action of $G_{+}$ exhibits the family as a
torsor under $G_{+}$, which the exponent identifies with $(\mathbb{R}, +)$.
\end{proof}

Theorem~\ref{thm:torsor} is the structural heart of the paper. It says the
remaining freedom in the cost is not a vague ``one parameter'' but a gauge: the
admissible costs form a single orbit, all equivalent under the symmetry group of
the domain, with no member algebraically preferred over any other. This is the
exact sense in which the cost \emph{form} is forced. The form is the orbit. What
is not forced is the choice of representative within the orbit, and that choice is
one real number.

\section{One datum fixes the gauge; calibration selects the canonical cost}

A torsor has no distinguished point until one is named. Naming one point requires
exactly one piece of data, and we now show that a single comparison value
suffices, after which calibration identifies the named point with the canonical
cost.

\begin{theorem}[A single datum pins the gauge]\label{thm:single}
Let $\costL{\lambda}$ and $\costL{\nu}$ be two members of the family, and suppose
they agree at a single ratio $x_0 > 1$, that is $\costL{\lambda}(x_0) = \costL{\nu}(x_0)$.
Then $\lambda = \nu$, so the two costs coincide everywhere. Equivalently, the
value of the cost at any one ratio different from the unit determines the gauge
coordinate.
\end{theorem}

\begin{proof}
Fix $x_0 > 1$, so $\ln x_0 > 0$. The map $\lambda \mapsto \costL{\lambda}(x_0)
= \cosh(\lambda \ln x_0) - 1$ is strictly increasing for $\lambda > 0$, because
$\cosh$ is strictly increasing on $(0, \infty)$ and the argument $\lambda \ln x_0$
is strictly increasing in $\lambda$. A strictly increasing function is injective,
so equal values at $x_0$ force $\lambda = \nu$.
\end{proof}

\begin{theorem}[Calibration selects the canonical cost]\label{thm:calib}
Among the family $\{\costL{\lambda} : \lambda > 0\}$, the unit calibration (L5)
holds if and only if $\lambda = 1$, and the calibrated member is the canonical
cost: $\costL{1} = J$.
\end{theorem}

\begin{proof}
In logarithmic coordinates the profile of $\costL{\lambda}$ is
$t \mapsto \cosh(\lambda t) - 1$, whose second derivative at $t = 0$ is
$\lambda^{2}$. Calibration requires this to equal $1$, that is $\lambda^{2} = 1$,
and for $\lambda > 0$ this holds exactly when $\lambda = 1$. At $\lambda = 1$,
$\costL{1}(x) = \tfrac{1}{2}(x + x^{-1}) - 1 = J(x)$.
\end{proof}

Putting the two together, the gauge is fixed by one datum, and the particular
datum ``the logarithmic curvature at the unit equals one'' is the choice that
names the canonical cost. The calibration is the unit speed of the gauge group:
in the exponent coordinate the curvature reads off as $\lambda^{2}$, so demanding
unit curvature is demanding $\lambda = 1$, the identity of the gauge group, which
is the canonical cost.

\section{The honest stratification}\label{sec:stratification}

We now assemble the discrete negative and the continuous positive into the single
stratified statement that is the conclusion of the paper.

\begin{theorem}[Cost stratification]\label{thm:strat}
The status of the canonical reciprocal cost $J$ under the cost laws is as follows.
\begin{enumerate}[leftmargin=2em, label=\textup{(\alph*)}]
  \item \textbf{Discrete underdetermination.} On the rational field that
    distinction generates, the laws (L1) through (L3) do not force $J$. The
    admissible cost forms are the character costs, an infinite product with one
    free factor per prime, and $J$ is one point of it
    \textup{(Theorem~\ref{thm:negative})}.
  \item \textbf{Continuous form forcing.} On the continuous completion, the laws
    (L1) through (L4) force the cost form to lie in the one-parameter family
    $\{\costL{\lambda} : \lambda > 0\}$, and this family is exactly the solution
    set \textup{(Theorem~\ref{thm:family})}.
  \item \textbf{The form is a gauge orbit.} That family is the orbit of $J$ under
    the gauge group of multiplicative automorphisms of the positive reals, acted
    on freely and transitively, hence a torsor; the residual freedom is exactly
    one real degree \textup{(Theorem~\ref{thm:torsor})}.
  \item \textbf{One datum, then calibration.} A single comparison value fixes the
    gauge \textup{(Theorem~\ref{thm:single})}, and the unit calibration (L5)
    fixes it to $J$ \textup{(Theorem~\ref{thm:calib})}.
\end{enumerate}
The honest reading is therefore: distinction and the algebraic laws force the cost
\emph{form}, the gauge orbit of $J$; a single unit of scale, supplied by one
calibration datum, selects $J$ within the form; and the value of that unit is a
multiplicative-automorphism gauge that distinction does not, by itself, fix.
\end{theorem}

This is weaker than ``distinction forces $J$,'' and it is true, where the stronger
claim is false. It is also more informative. It says precisely what is forced, the
form; precisely what remains free, one real gauge coordinate; and precisely what
fixes the freedom, one datum. The total arbitrary content of the cost layer is
exactly one positive real number. Everything else about the cost, its reciprocal
symmetry, its composition behavior, the cosh shape of its profile, the fact that
it is the orbit of a single canonical representative, is forced.

\section{The one open question that would upgrade the result}\label{sec:open}

There is a single question whose answer would change the stratification of
Theorem~\ref{thm:strat} into a clean forcing, and we name it precisely so that it
is not mistaken for something already settled.

The one free datum is the unit calibration, the demand that the logarithmic
curvature of the cost at the unit equals one. In the present account this is
supplied from outside, as a choice of scale. The open question is whether that
curvature is itself derivable from distinction, specifically from the cost of a
single distinction act. The natural conjecture is that the minimal recognition
cost of telling apart two states that differ by one act, carried from the
discrete layer to the completion, fixes the second-order coefficient of the cost
profile at the unit to exactly one. If that conjecture is a theorem, then the
calibration datum is not a free choice but a consequence of the primitive, the
continuous classification runs with no free input, and ``distinction forces $J$ on
the completion'' becomes true without circularity. If the conjecture is false,
then the genuine native curvature is some other value, the forced cost is the
corresponding member of the family rather than $J$, and one has learned the true
cost of distinction, which is itself a result.

We do not claim either outcome here. The point of stating the question is to mark
the exact location of the remaining freedom. It is not spread across the cost
function; it is concentrated in one number, the curvature at the unit, and the
question of whether distinction fixes that number is a single, sharp, falsifiable
target. Until it is settled, the honest claim is the stratified one of
Theorem~\ref{thm:strat}: the form is forced, and the unit is a gauge.

\section{Conclusion}

The load-bearing joint of the program is the cost function, and we have measured
exactly how far it is forced. On the discrete field of ratios that distinction
generates, the cost is not forced at all: the admissible costs are the character
costs of a free abelian group, an infinite product with one independent factor
per prime, and the canonical cost is one unremarkable point. Any uniqueness on the
discrete field must assume the cost on a generating set, which is circular as a
forcing claim. The freedom is removed only by passing to the continuous
completion, where continuity restricts characters to power maps and collapses the
infinite family to a single one-parameter family. That family is exactly the orbit
of the canonical cost under the multiplicative automorphisms of the line, a torsor
with one real degree of freedom, and a single calibration datum names the
canonical representative within it.

The result is a stratification, not a forcing, and the stratification is the
honest and the stronger statement. It localizes the entire arbitrary content of
the cost layer to one real number, the unit of scale, and proves everything else
forced. It identifies the single open question, whether that one number is itself
the cost of a distinction act, that would close the gap entirely. And it does all
of this without the circularity that a naive forcing claim on the rationals would
have concealed. The cost form is forced. The unit is a gauge. Naming which is
which, with proofs on both sides, is what it means to take the joint seriously.

\begin{thebibliography}{99}

\bibitem{Aczel1966}
J.~Acz\'el,
\emph{Lectures on Functional Equations and Their Applications},
Academic Press, 1966.

\bibitem{Aczel1989}
J.~Acz\'el and J.~Dhombres,
\emph{Functional Equations in Several Variables},
Cambridge University Press, 1989.

\bibitem{dAlembert1769}
J.~le~Rond~d'Alembert,
\emph{Opuscules math\'ematiques}, vol.~6, 1769.

\bibitem{WashburnFE2026}
J.~Washburn,
``A Functional-Equation Encoding of Logical Consistency on Continuous
Positive-Ratio Comparisons: A Conditional Rigidity Theorem,'' draft, 2026.

\bibitem{WashburnContinuum2026}
J.~Washburn,
``Where Necessity Ends: Why Distinction Forces Arithmetic but Not the
Continuum,'' draft, 2026.

\bibitem{WashburnArithmetic2026}
J.~Washburn,
``Arithmetic from Distinction: A Number Tower Forced by a Single Primitive,''
draft, 2026.

\bibitem{WashburnZlatanovic2026}
J.~Washburn and Z.~Zlatanovi\'c,
``Uniqueness of the Canonical Reciprocal Cost,''
\emph{Axioms}, 2026.

\end{thebibliography}

\end{document}
